% !TeX root = ../../Book3.tex
\section{调和测度}
\label{b3:sec:1604}

正则域上的解算子保持线性与正性。测度表示把整个解算子压缩成一族边界概率测度，共形映射再将这些测度从一个区域搬到另一个区域。

\subsection{解的表示}
\label{b3:subsec:160401}

设 \(D\) 是有界正则域。固定 \(z\in D\)，映射 \(\varphi\mapsto H_\varphi(z)\) 为 \(C(\partial D)\) 上的正线性泛函：两解相加仍给相加边值，唯一性使解算子线性；非负边值由最大值原理给非负解。常数一的解为一。紧空间上的正泛函表示遂给唯一概率测度 \(\omega_D^z\)，满足
\[
H_\varphi(z)=\int_{\partial D}\varphi\,\mathrm{d}\omega_D^z.
\]
这里使用第一卷定理11.4.1的 Riesz--Markov 表示。对于不同 \(z\)，这些测度通常互不相同；Harnack 估计使其对非负连续测试函数的积分在内部紧集上相互控制。

\subsection{共形不变性}
\label{b3:subsec:160402}

\begin{proposition}\label{b3:prop:16-harmonic-measure-conformal}
若 \(F:D\to D'\) 共形双射，且延拓为闭包间同胚，两域均有界正则，则
\(F_\#\omega_D^z=\omega_{D'}^{F(z)}\)。
\end{proposition}
\begin{proof}
对 \(\varphi\in C(\partial D')\)，\(H_\varphi\circ F\) 调和，因为 \(\Delta(u\circ F)=|F'|^2\cdot(\Delta u)\circ F\)。它连续取边值 \(\varphi\circ F\)，故由唯一性等于 \(H_{\varphi\circ F}\)。两边在 \(z\) 的值给两测度对全部连续测试函数积分相等，表示的唯一性得结论。
\end{proof}
若映射只在内部存在，边界像可能需要素端而不是通常拓扑边界。此时不能无说明地写相同的推前公式。

\subsection{概率解释的预告}
\label{b3:subsec:160403}

概率论中，调和测度可解释为平面 Brownian 运动首次退出区域时的位置分布。这个解释需要先构造过程，再证明退出时间、停时和边界到达的性质；它不是概率测度表示定理本身的内容。

圆盘中心的退出分布由旋转对称性应当均匀，与 Poisson 核的中心值一致。非中心处则偏向附近边界弧。本章使用的是已构造的解与测度，不以这一直观概率解释代替 Perron 证明。


\ProseParagraphBreak

解表示蕴含测度之间的定量关系。取 \(z,w\) 属于同一内部紧集，用有限 Harnack 链得到 \(C^{-1}H_\varphi(w)\le H_\varphi(z)\le C H_\varphi(w)\) 对所有非负连续 \(\varphi\) 成立。正测度的序关系因此给
\[
C^{-1}\omega_D^w(E)\le\omega_D^z(E)\le C\omega_D^w(E)
\]
对每个 Borel 边界集合 \(E\) 成立：先由连续函数积分不等式得到测度支配，再用正则性从紧集、开集推广。尤其是否具有零调和测度与内部观察点无关，虽然具体权重随点而变。这个结论只比较同一区域内的点；跨域比较需要域的关系或共形映射。

\ProseParagraphBreak

对边界弧 \(I=[\alpha,\beta]\) 的圆盘权重，中心值为 \((\beta-\alpha)/(2\pi)\)。非中心处可积分核，或用 Möbius 变换 \(T_z(\zeta)=(\zeta-z)/(1-\overline z\zeta)\) 将 \(z\) 搬到原点，得到 \(\omega^z(I)=|T_z(I)|/(2\pi)\)，这里弧长在像的单位圆上测量。原来的近弧会被展开，远弧会被压缩；这把 Poisson 核分母 \(|e^{\mathrm{i}t}-z|^2\) 中的距离效应转成圆周变换的导数。必须取相应连续像弧，不能只减两个主辐角而忽略跨越分支切线的情形。

\ProseParagraphBreak

例如上半平面的区间 \([a,b]\) 在 \(x+\mathrm{i}y\) 处的权重是
\[
\omega^{x+\mathrm{i}y}([a,b])=\frac1\pi\left(\arctan\frac{b-x}{y}-\arctan\frac{a-x}{y}\right).
\]
取 \(y\downarrow0\)，当 \(x\in(a,b)\) 时极限一，区间外极限零，端点径向极限二分之一。这个算式既检验归一化，也显示端点与弧内部的边界行为不同。若同时改变 \(x\) 和 \(y\) 以斜向接近端点，比例 \((a-x)/y\) 可趋于不同实数，因而极限也相应变化。

\ProseParagraphBreak

调和测度一般不等于普通边界长度；即使圆盘内也只有中心处是均匀的。对有尖角的域，共形变换在边界的伸缩会加强某些位置的权重，正则域这一性质本身没有承诺与弧长的密度存在。下一章光滑边界的 Green 函数法向导数将给具体密度公式，其额外边界光滑性正是微分边界核所需的条件。
